\section{Introduction}
\label{sec:intro}

Security teams are starting to hand alert triage and investigation to LLM agents~\citep{guide,cortex,sentinelrl}. An agent receives an alert, queries logs, and reports whether an attack happened and what it touched. Whether such an agent can be trusted to close alerts is decided largely by benchmarks: ExCyTIn-Bench~\citep{excytin}, SIABench~\citep{siabench}, and others report how often agents reach the right answer, and the field reads the numbers as measures of investigation ability. Investigation is the reason to deploy an agent at all. An agent that only recognizes what the alert already says adds nothing an analyst would pay for, and it fails in the cases an attacker creates on purpose, where the malicious activity looks routine and nothing names it.

A benchmark measures investigation only if its items cannot be answered without investigating. Consider an ExCyTIn question generated from a three-hop path: from an alert about a scheduled task, through the host, to a related alert whose IP address is the answer. The question asks for ``the IP address recorded in the related `Anomaly detected in ASEP registry' alert''. One SQL query for that alert name returns the answer; the pivot the question was generated from is unnecessary (\cref{fig:example}). If many items are like this, an agent can earn a score without investigating, and the score says nothing about the cases that matter.

Benchmark papers report model scores but rarely the score of a rule that does not investigate. Security has a tradition of re-evaluations that ask exactly this question of learning-based detectors: simple baselines match complex provenance-based detectors~\citep{bilot2025simpler}, and spurious correlations inflate lab results~\citep{arp2022dos}. Machine learning knows the same problem as shortcut learning~\citep{geirhos2020shortcut,gururangan2018artifacts}. Recent work on incident response proposes new benchmarks that reward evidence discovery~\citep{sirbench}, but the benchmarks already in use have not been audited, and whether the agents evaluated on them use the shortcuts they contain has not been measured.

\begin{figure}[t]
  \centering
  \includegraphics[width=\columnwidth]{fig_floor.pdf}
  \caption{A lookup without an LLM outscores \oBelowGraph{} of \oModels{} reported agents on ExCyTIn-Bench. Bars: mean reward reported for each model on the 589 test questions~\citep{excytin}. Lines: two investigation-free lookups of the alert the question names, on the same questions, scored by exact match (shaded: 95\% intervals, incidents resampled). Dark bars exceed the graph lookup.}
  \label{fig:floor}
\end{figure}

We audit four public benchmarks and datasets for LLM security-operations agents (ExCyTIn-Bench, SIABench, GUIDE, and OTRF Security-Datasets) and three triage environments we built ourselves. For each, we write a fixed rule that uses only what the agent sees, contains no LLM, and skips the investigation the benchmark intends, and we compare it with reported model scores. We find four kinds of shortcut: the question names the evidence that holds the answer (S1), the label follows the organization that assigned it (S2), the data generation leaves a cue such as the attack network's address (S3), and in constructed environments a fixed procedure over the evidence model solves every case (S4). For ExCyTIn, whose authors publish the full logs of 14 agent runs, we go further and test question by question whether the agents' successes come from the shortcut. We registered the hypotheses and froze the rules before every confirmatory run.

\paragraph{Results} Every audited benchmark admits an investigation-free rule that reaches a large share of the reported scores. On ExCyTIn, \oNamed\% of the test questions name the alert that holds the answer, and a lookup of that alert answers \oGraph\% exactly, more than the mean reward of \oBelowGraph{} of \oModels{} reported models (\cref{fig:floor}). On SIABench, ``true positive if and only if the alert involves 172.16.0.1'' is correct on all 35 alerts. On GUIDE, repeating an organization's earlier grade for the same detector is correct on \gHistPct\% of later incidents. In OTRF recordings, the alert text alone separates \otRuleCorrect{} of \otN{} alerts. Yet the ExCyTIn agents do not use the shortcut. Their success is no higher on questions that name the answer's alert (\lgPos{} of \lgModels{} runs, $p=\lgP$); they issue as many queries when one would do; and their ranking is nearly unchanged on shortcut-resistant questions (Kendall $\tau=\lgTau$). The deeper problem is that success depends only weakly on the investigation a question requires: the strongest agent fails \tcOpusFail{} of the 140 questions a single lookup answers.

\paragraph{Contributions}
\begin{itemize}\setlength\itemsep{1pt}
\item A definition of shortcuts for investigation benchmarks and a taxonomy of four kinds, each paired with an investigation-free baseline (\cref{sec:method}).
\item An audit of four public benchmarks and three constructed environments, showing that each admits a rule that reaches or exceeds a large part of reported model scores (\cref{sec:floors}).
\item A question-level analysis of 14 published agent runs on ExCyTIn, showing that the agents do not exploit the shortcut and that their success tracks the required investigation only weakly (\cref{sec:scores}).
\item Checks for benchmark builders, evidence that paired attack and benign ``twins'' remove the single-observation shortcut, and per-question shortcut flags for ExCyTIn (\cref{sec:design}).
\end{itemize}
